\section{Quarter-turn refinement and octahedral sectors}
\label{sec:quarter}
\begin{theorem}\label{thm:quarter}
Assume the pinned spatial-transfer construction of \cite{upstream}, the
recorded untwisted and $P$-twisted thermal inputs, the cyclic inputs of
version 1.6, the retained 120-site trial, and the quarter-turn and
84-site integer certificates described below. At coupling $J=1$,
\[
 \boxed{\gap_L>\frac{303}{2000}=0.1515
 \quad\text{for every integer }L\ge24.}
\]
The ground state is unique. The strict lower bounds $0.063$, $0.119$,
and $0.145$ hold for every integer length at least $18$, $20$, and $22$,
respectively. Separately, $\gap_{18}>0.088$.
In particular $\liminf_{L\to\infty}\gap_L\ge0.1515$ over all integer
lengths, including odd lengths.
\end{theorem}
The new input is a fourth boundary rotation, through $\pi/2$. Its
thermal traces resolve the octahedral representations more finely.
The resulting sector identities are explicit consequences of the
upstream spatial representation and physical rotational covariance;
no new general representation-theoretic principle is claimed.
The scalar recurrence and tight purity extraction are retained from
versions 1.5 and 1.6. Earlier proofs and records remain available.

\subsection{Four resolved lists}
The proper signed permutation group $\mathcal G$ in
\cite[Proposition 3.1]{upstream} has 24 elements and is isomorphic to
$S_4$ through its action on the four unoriented body diagonals of a cube.
Let $P_f$ be a coordinate-axis half-turn, $P_e$ an edge-axis half-turn,
$C$ a third-turn, and $Q$ a coordinate-axis quarter-turn. The character
table, in the order used here, is
\begin{center}
\begin{tabular}{c|rrrrr}
 & $1$ & $P_f$ & $P_e$ & $C$ & $Q$\\
Class size & 1 & 3 & 6 & 8 & 6\\\hline
$A$ & 1 & 1 & 1 & 1 & 1\\
$B$ & 1 & 1 & -1 & 1 & -1\\
$E$ & 2 & 2 & 0 & -1 & 0\\
$T$ & 3 & -1 & 1 & 0 & -1\\
$V$ & 3 & -1 & -1 & 0 & 1
\end{tabular}
\end{center}
These representations can be constructed as the trivial and sign
representations, the two-dimensional standard representation on the
three pairings, the standard representation on four diagonals, and
its sign twist. The last has the character of the natural rotation
representation. The finite checker enumerates the rotations, verifies
character orthogonality, and checks all central-projector products
\[
 \Pi_\rho=\frac{d_\rho}{24}\sum_{g\in\mathcal G}
                  \chi_\rho(g)U_g.
\]
This verifies the finite group algebra; the representation on the
spatial Hilbert space remains the cited analytic premise.

Because $X_b$ commutes with the group, its spectrum decomposes into
real multiplicity-space lists $A,B,E,T,V$ with representation
multiplicities $1,1,2,3,3$. Their moments converge absolutely for orders
at least two. Subscripts below indicate moment order, not a physical
spin quantum number.

\begin{lemma}\label{lem:octahedral}
For every $b>0$ and integer $n\ge4$, the sector moments satisfy
$T_n=B_n+E_n$, and
\begin{align*}
 Z_n&=A_n+4B_n+5E_n+3V_n,\\
 Z_n(P)&=A_n+E_n-V_n,\\
 Z_n(C)&=A_n+B_n-E_n,\\
 Z_n(Q)&=A_n-2B_n-E_n+V_n.
\end{align*}
Consequently
\begin{align*}
 A_n&=[Z_n+9Z_n(P)+8Z_n(C)+6Z_n(Q)]/24,\\
 B_n&=[Z_n-3Z_n(P)+8Z_n(C)-6Z_n(Q)]/24,\\
 E_n&=[Z_n+3Z_n(P)-4Z_n(C)]/12,\\
 V_n&=[Z_n-3Z_n(P)+2Z_n(Q)]/8.
\end{align*}
The former lists obey $I=A+B$, $O=E$, and $N=B+E+V$,
with equality of nonzero eigenvalue multiplicities.
\end{lemma}
\begin{proof}
Both half-turns are conjugate in $SO(3)$, so a simultaneous physical
rotation at every site gives $Z_n(b,P_e)=Z_n(b,P_f)$.
They need not be conjugate inside $\mathcal G$.
Subtracting their character columns in the spatial trace identity
gives $-2B_n-2E_n+2T_n=0$. Substitution in the other columns gives the
four displayed traces, and direct inversion gives the sector formulas.
Restriction to the subgroup generated by $D_2$ and $C$ gives
$I=A+B$, $O=E$, and $N=T+V$.

For completeness, the moment relation also identifies multiplicities.
On a compact interval containing the spectrum, form the finite signed
measure with atom at each nonzero $\mu$ of weight
$\mu^4(m_T(\mu)-m_B(\mu)-m_E(\mu))$.
Absolute summability follows from the fourth moments. All its polynomial
moments vanish by the identity at orders $n\ge4$.
Polynomial density in continuous functions on that compact interval
forces the measure to vanish. Thus the multiplicities agree at every
nonzero eigenvalue. In the full spectrum the four remaining lists
therefore have effective multiplicities $1,4,5,3$.
\end{proof}
No statement about the degeneracy of the first \emph{physical}
excitation follows from these spatial multiplicities.

\subsection{Integer quarter-turn inputs}
All nine traces $Z_n(49/4,Q)$, $4\le n\le12$, are computed by
\path{compute_quarter_moments.cpp}, an extension of the version 1.6
integer Horner evaluator. For the new twist the seam coefficients are
$\pm i$. Coordinates are Gaussian integers $u+iv$, and the final squared
column norm is $\sum(u^2+v^2)$.
The divisor 32, precision $Q_*=2^{55}$, cutoff 280, Poisson mean 98,
and exact coefficient floors are unchanged.

The local energy argument in Lemma~\ref{lem:local-energy} applies to
$Q$ as well: every open five-site path can have its seam removed by
on-site unitary rotations. Hence the same contraction modulus
$\kappa_n=1-n/384$ is valid for $5\le n\le12$; use one at $n=4$.
In a block of dimension $d$, two coordinate floors now have vector
error at most $\sqrt{2d}$. Set $d^+=\lfloor\sqrt{2d}\rfloor+1$ and use
\[
 e_{n,d}=\begin{cases}
 281(d^++1)+1,&n=4,\\
 \lceil384(d^++1)/n\rceil+1,&5\le n\le12.
 \end{cases}
\]
As before, $E_n=\lfloor\sqrt{\sum_M d_Me_{n,d_M}^2}\rfloor+1$
bounds the Frobenius error in $Q_*$ units. The final unit includes the
same conditioned Poisson-tail estimate. The representative-column
symmetries remain valid for this unit-modulus seam phase. Bounds on
coordinate row sums and column norms verify that intermediate signed
64-bit and 128-bit integers, and the checked 256-bit accumulator,
cannot overflow.

The complete table was recomputed from the distributed C++ source.
Independent all-column Gaussian-integer calculations at lengths
4, 5, and 6, without orbit reduction, give overlapping rigorous norm
intervals. Untwisted and half-turn four-site totals also match the
unmodified upstream packed implementation exactly. These are internal
implementation checks, not external review.

\subsection{Fixed moment certificates and initialization}
Applying eight fixed filters $x^4p(x)^2$, with quartic $p$,
to the resolved moments gives
\begin{align*}
 -.5731783&<\mu_A<1.0007981,\\
 -.5319629&<\mu_B<.4316513,\\
 -.5775075&<\mu_E<.6993004,\\
 -.8668243&<\mu_V<.7170568.
\end{align*}
Every filter is nonnegative on the real line, and its shifted
coefficients certify the whole exterior ray. The old $N$ list now lies
in $(-.8668243,.7170568)$ by Lemma~\ref{lem:octahedral}.
This is useful for short odd lengths, where signed contributions matter.

The certificate verifies 117 fixed rational polynomial majorants and
minorants on these complete intervals, using exact Bernstein subdivision.
They combine moment orders 4 through 12. The coefficient records are in
\path{independent_results/quarter_moment_polynomials.json}.
Capped-mass bounds are retained when smaller. Reconstructing
$I=A+B$, $O=E$, and $N=B+E+V$ then supplies the inputs to the retained
three-sector recurrence; no numerical optimizer enters acceptance.

The starting choice is
\[
 m=26,\qquad t=2,\qquad n_0=42,\qquad b_0=49/2.
\]
The published integer trial matrices are newly contracted at length 84.
Their norm $V_{84}$ and one-bond numerator $U_{84}$ satisfy
\[
 V_{84}>0,\qquad500000U_{84}+700741V_{84}<0.
\]
A four-site direct wavefunction checks the contraction convention.
Thus $Z_{84}(b)>1$ for every $b>0$, giving the physical-purity seed as
before. The retained 120-site trial gives the spatial leading-eigenvalue
bound. The resulting upper defects and residuals are approximately
\[
 i=.008031531,\quad j=.008876122,\quad
 r=.037366075,\quad q=.009544010.
\]
These displays are rounded upward; the certificate stores the exact
accepted fractions.

\subsection{Complete length cover and scope}
The new checker has 381 named exact checks, in addition to root-direction
and full-interval polynomial comparisons. It checks every length
24 through 62 individually, followed by these seven intervals:
\begin{center}
\begin{tabular}{rrr}
\toprule First length & Last length & Recurrence level\\\midrule
63 & 65 & 0\\
66 & 115 & 0\\
116 & 328 & 0\\
329 & 1370 & 2\\
1371 & 5831 & 4\\
5832 & 29770 & 6\\
29771 & 32255 & 7\\\bottomrule
\end{tabular}
\end{center}
At reference length 21504 and inverse temperature 12544, the retained
invariant starts with the exact rational
$w=1.70982022810726\times10^{-827}$, also recorded in
\path{independent_results/quarter_refinement.json}. It covers every
integer length from 32256 onward. The checker verifies its starting
inequalities, propagation with $w_+=6w^2$, and strict exponential
comparisons for the whole infinite tail. The shorter-chain bounds in
Theorem~\ref{thm:quarter} are checked separately and joined to the
stronger ranges; no monotonicity of the gap in chain length is assumed.

The quarter-turn data and the refined moment decomposition produce the
gain. An exploratory recurrence retaining additional octahedral
projection caps gave no material further improvement in the screened
parameters; it is not used in the proof. The original purity-only
squaring estimate is saturated by a two-level probability distribution,
so it cannot be improved universally without further information.
These observations do not prove optimality of the present method or
parameters. The result remains conditional on the upstream analytic
construction, unreviewed, and not machine-formalized.
